\section{灵敏度与稳健性检验}

\subsection{检验设计}

灵敏度文件包含五类局部证据：导引模型、丢失导引反事实、烟幕漂移、响应延迟和 UAV 作用半径。每一行包含类别、参数、参数值、指标和结果；正式基线与实验性丢失导引反事实在 JSON 中分开标记。该设计是局部扰动，不是针对所有参数的全局不确定性分析。

对于固定参考候选，漂移敏感度采用离散检测时间网格 $\{t_\ell\}$，计算最小间隙
\begin{equation}
  g_{\min}(\delta)=\min_{\ell}\left(
  \|x_T(t_\ell)-[x_S+\Delta(\delta)]\|_2
  -r_T-R_S(t_\ell)\right),
  \label{eq:sensitivity-gap}
\end{equation}
其中 $\Delta(\delta)$ 是大小为 $\delta$ 的有界漂移。响应延迟测试只改变最早合法爆炸时刻，因此可用窗口为
\begin{equation}
  W(\tau)=\max\left(0,\ t_{\mathrm{det,end}}-
  \max(t_{\mathrm{det,start}},\tau+\tau_{\mathrm{burst}})\right).
  \label{eq:sensitivity-window}
\end{equation}

\begin{table}[!htbp]
\centering
\caption{灵敏度结果中的代表性边界}
\label{tab:sensitivity}
\begin{tabular}{llrr}
\toprule
类别 & 参数/指标 & 低值结果 & 高值结果 \\
\midrule
导引模型 & 检测时长/s & 24.1677 & 24.1677 \\
烟幕漂移 & 最小间隙/m（漂移 0--40 m） & -39.6683 & 0.0014 \\
响应延迟 & 可用窗口/s（延迟 1.5--3.0 s） & 19.1677 & 17.6677 \\
作用半径 & 路径可行性（8000--14000 m） & 真 & 真 \\
丢失导引 & 最小分离/m（实验反事实） & 69.4180 & 69.4180 \\
\bottomrule
\end{tabular}
\end{table}

表\ref{tab:sensitivity}的导引模型结果来自瞬时纯追踪与惯性航向响应率 $0.5$、$1.0$、$2.0\,\mathrm{s}^{-1}$ 的检测时长；比例导航项因缺少导航比和过载来源而在结果文件中标记为 deferred。烟幕漂移的最小间隙从负值接近零，说明固定中心方案对该扰动较敏感；响应延迟增加会压缩可用检测窗口。

\begin{figure}[!htbp]
  \centering
  \includegraphics[width=0.86\textwidth]{sensitivity_summary.png}
  \caption{冻结基线中的 Stage 6 局部灵敏度汇总图。}
  \label{fig:sensitivity}
\end{figure}

\subsection{正式与反事实的隔离}

丢失导引参数 $(\tau_T,\tau_L,T_R)$ 的 $27$ 行实验结果只用于说明失锁模型变化下的几何分离，不修改正式 Q1 候选决策，也不改变 Q2--Q4 任务包。正式基线字段为 \texttt{formal\_baseline=true}，反事实对象为 \texttt{formal\_baseline=false}。因此，本文不把反事实数值写成正式方案的稳健保证。

\subsection{稳健性结论}

在代码已计算的局部范围内，UAV 作用半径从 $8000$ 增至 $14000\,\mathrm{m}$ 时固定路径均被判定为可行；响应延迟每增加 $0.5\,\mathrm{s}$，可用窗口减少约 $0.5\,\mathrm{s}$。这些是有限样本、固定路径或固定候选下的敏感性方向，不足以支持完整参数空间中的稳健最优性结论。
